\section*{Chapitre \ref{ch:s1:options-implied-signals-for-stocks} --- Signaux implicites des options sur les actions}

\begin{solution}{exo:s1:options-implied-signals-for-stocks:1}
Une mauvaise nouvelle est pondérée par 1.5~: $0.003 \times 1.5 \times 2 = 0.009$, soit 0.9 point de volatilité sur les puts~; une bonne nouvelle donne $0.003 \times 2 = 0.006$, soit 0.6 point sur les calls. Les deux
sont inférieurs au bruit quotidien d'un point, ce qui explique pourquoi une journée isolée n'est pas informative et que le signal doit être lu par rapport à son historique.
\end{solution}

\begin{solution}{exo:s1:options-implied-signals-for-stocks:2}
$0.34~\% \times 52 = 17.7~\%$ par an et $0.5~\% \times 52 = 26~\%$ par an, avant coûts et sans capitalisation.
\end{solution}

\begin{solution}{exo:s1:options-implied-signals-for-stocks:3}
Un skew prononcé signifie que les puts sont poussés à la hausse, ce qui est le cas lorsque les traders informés anticipent une baisse~: l'action a plus de chances de baisser, donc le skew entre avec un signe moins. Un
écart call--put élevé signifie que les calls sont poussés à la hausse par rapport aux puts, ce qui est le cas lorsque de bonnes nouvelles sont attendues~: il entre avec un signe plus.
\end{solution}

\begin{solution}{exo:s1:options-implied-signals-for-stocks:4}
Vendre l'action à découvert coûte des frais d'emprunt, exige un locate et peut être impossible ou donner lieu à un rappel~; acheter un put n'entraîne aucun de ces coûts et ajoute de l'effet de levier. Pour une bonne
nouvelle, acheter l'action est peu coûteux et facile, si bien que les options n'offrent que de l'effet de levier.
\end{solution}

\begin{solution}{exo:s1:options-implied-signals-for-stocks:5}
Les traders informés n'agissent qu'avant les événements~: les signaux ne contiennent donc de l'information qu'à ces moments-là~; ailleurs, ce sont du bruit qui dilue l'IC. Le livre doit
négocier les signaux sur les actions dont un événement approche et rester à l'écart des autres.
\end{solution}

\begin{solution}{exo:s1:options-implied-signals-for-stocks:6}
Les volatilités implicites portent un bruit quotidien d'un point, si bien que le classement hebdomadaire change beaucoup d'une semaine à l'autre~: le livre négocie beaucoup pour un petit
signal diffus, et dix points de base par unité négociée coûtent davantage que les 8.5~\% par an qu'il rapporte brut.
\end{solution}

\begin{solution}{exo:s1:options-implied-signals-for-stocks:7}
Le livre sur les volumes autour des événements rapporte 9.1~\% par an après coûts (ratio de Sharpe 1.75), contre 23.0~\% et 31.9~\% pour le skew et l'écart. Le volume augmente avec la taille de
la demande informée, quel qu'en soit le sens~; seule la pondération supplémentaire des mauvaises nouvelles lui donne un signe, si bien que son IC vaut à peu près la moitié de celui des signaux de prix.
\end{solution}

\begin{solution}{exo:s1:options-implied-signals-for-stocks:8}
Le chiffre provient d'un modèle dans lequel les traders informés connaissent la surprise dans son intégralité et où leur demande en est une fonction simple. Le flux informé réel est faible,
non signé dans la plupart des données, mêlé à la couverture et disputé~; les effets publiés se chiffrent en dizaines de points de base par semaine avant coûts. Il faut s'attendre à une
petite fraction, et tester sur des données réelles d'options, coûts compris.
\end{solution}

\begin{solution}{pb:s1:options-implied-signals-for-stocks:1}
\begin{enumerate}
\item \omterm{def:s1:options-implied-signals-for-stocks:informed}{Trading d'options informé} par les traders disposant d'une information privée~; volatilité implicite du call moins celle du put à prix d'exercice appariés~; volume des puts rapporté au volume des calls~; \omterm{def:s1:options-implied-signals-for-stocks:volume}{ratio de volume options
sur actions}.
\item Les skews les plus marqués sous-performent les plus plats de 10.9~\% par an ajustés du risque, pendant au moins six mois, avec des chocs de résultats plus défavorables au trimestre suivant.
\item Les calls chers surpassent les puts chers de 50 points de base par semaine, un écart qui s'amenuise avec le temps~; l'écart call--put a un pouvoir prédictif positif, la volatilité réalisée moins l'implicite
a un pouvoir prédictif négatif.
\item Les faibles \omterm{def:s1:options-implied-signals-for-stocks:volume}{ratios put--call} surpassent les forts de 40 points de base par jour et de plus de 1~\% par semaine~; un faible \omterm{def:s1:options-implied-signals-for-stocks:volume}{ratio de volume options sur actions} surpasse un fort de 0.34~\% par semaine.
\item Volatilité à la monnaie égale à la volatilité réalisée plus deux points, skew de put de trois points, un point de bruit quotidien~; les traders informés achètent des puts ou des calls dans les
dix jours avant les annonces, soit 0.3 point par écart-type de surprise.
\item Les coûts de vente à découvert rendent les options plus attractives pour les mauvaises nouvelles.
\item Il ne contient aucune information insérée à dessein, son IC devrait donc être nul, et il l'est.
\item Pour marquer les dix jours précédant chaque annonce, moment où les traders informés agissent.
\item Skew 0.004, 0.014, 0.018, 0.012~; \omterm{def:s1:options-implied-signals-for-stocks:spread}{écart de volatilité implicite} 0.005, 0.017, 0.022, 0.015~; volume 0.002, 0.008, 0.010, 0.007~; volatilité implicite moins réalisée nulle.
\item Skew 0.026, 0.085, 0.111, 0.082~; \omterm{def:s1:options-implied-signals-for-stocks:spread}{écart de volatilité implicite} 0.030, 0.101, 0.129, 0.094~; volume 0.015, 0.047, 0.061, 0.045.
\item Le saut lors de l'annonce survient dans les vingt jours suivant les achats informés~; au-delà, la fenêtre ajoute du bruit.
\item Toutes les actions~: ratios de Sharpe bruts de 3.17--3.67, négatifs après coûts~; autour des événements~: 23.0~\%, 31.9~\% et 9.1~\% par an après coûts pour le skew, l'écart et le volume.
\item \textbf{Résultat nommé.} L'IC de l'\omterm{def:s1:options-implied-signals-for-stocks:spread}{écart de volatilité implicite} vaut 0.005, 0.017, 0.022 et 0.015 à 1, 5, 20 et 60 jours pour l'ensemble des actions, et 0.030, 0.101, 0.129 et
0.094 parmi les actions annonçant dans les dix jours~; celui du skew~: 0.004, 0.014, 0.018, 0.012 et 0.026, 0.085, 0.111, 0.082.
\item Les traders informés du modèle connaissent la surprise avec exactitude et constituent la seule source de demande informée.
\item Faire du flux informé une faible part du volume d'options, ajouter les flux de couverture et de produits structurés, et laisser les traders informés choisir l'action lorsqu'elle est moins chère.
\item Lisser les signaux, ne les négocier que lorsque les événements sont proches et rééquilibrer moins souvent.
\item Surfaces de volatilité implicite par action, volume d'options (signé si possible), calendrier des événements et frais d'emprunt.
\item L'écart call--put (Cremers et Weinbaum).
\item Ils négocient les mêmes événements~: les signaux tirés des options anticipent la surprise que la stratégie de dérive du chapitre~7 négocie après coup.
\item Parce que les traders informés négocient souvent les options en premier.
\end{enumerate}
\end{solution}

\begin{solution}{iq:s1:options-implied-signals-for-stocks:1}
Les traders informés utilisent les options pour l'effet de levier et pour éviter les coûts de vente à découvert~; leur demande fait bouger les prix et les volumes d'options avant que l'action ne bouge, et le marché des
actions est lent à suivre.
\end{solution}

\begin{solution}{iq:s1:options-implied-signals-for-stocks:2}
À partir d'une surface ajustée à une maturité fixe (par exemple 30 jours)~: la volatilité implicite pour un moneyness fixe hors de la monnaie (un delta de $-0.2$ ou un strike à
90~\%) moins la volatilité implicite à la monnaie, à partir des cotations de clôture, comparée à l'historique propre de l'action et au skew du marché.
\end{solution}

\begin{solution}{iq:s1:options-implied-signals-for-stocks:3}
L'écart call--put compare deux options sur la même action et mesure la demande relative pour la hausse face à la baisse, un signal d'information~; la volatilité implicite moins
réalisée compare les options à l'historique propre de l'action et mesure la prime de risque de volatilité, un signal de risque.
\end{solution}

\begin{solution}{iq:s1:options-implied-signals-for-stocks:4}
Il faut vérifier si la transaction est à l'initiative de l'acheteur et si elle ouvre une position, la comparer au volume habituel avant publication des résultats, au skew et à l'encours vendeur à découvert de l'action~; si elle ressemble
à un achat informé de protection, il faut réduire ou couvrir les positions longues et considérer l'action comme sous-pondérée à l'approche de l'événement.
\end{solution}

\begin{solution}{iq:s1:options-implied-signals-for-stocks:5}
Cotations périmées ou non synchrones entre options et action~; rebond bid-ask dans les volatilités implicites~; primes d'exercice anticipé~; biais de survie dans les données d'options~; et
coûts de négociation de l'action sur des signaux bruités.
\end{solution}

\begin{solution}{iq:s1:options-implied-signals-for-stocks:6}
Avec une information privée d'une baisse d'ampleur $x$, la vente à découvert rapporte $x - c$ par unité~; un put à effet de levier $L$ rapporte environ $Lx$ moins la valeur temps de sa prime. Les
traders informés préfèrent les puts lorsque $Lx$ net de la prime dépasse $x - c$, ce qui est d'autant plus probable que $c$ est élevé~: lorsque la vente à découvert est coûteuse, les mauvaises nouvelles se reportent sur les puts, si bien que
le volume des puts (et le \omterm{def:s1:options-implied-signals-for-stocks:volume}{ratio de volume options sur actions}) porte davantage d'information négative que le volume des calls n'en porte de positive.
\end{solution}
